\chapter{Conclusion \& Discussion}
\label{chapter:conclusion_discussion}

This thesis studied a practical question for the next generation of personal and organisational agents: how can agents coordinate across people, teams, and organisations without turning every useful interaction into a privacy risk? The central claim is that cross-boundary delegation is not only a model-behaviour problem. It is an infrastructure problem. Once an agent holds private state, calls tools, remembers relationships, and talks to other agents, privacy depends on what the platform loads, what tools it exposes, and where it places the decision boundary.

The work contributes a setting, two benchmarks, and a set of architectural directions. SharedOS formalises cross-boundary delegation as a tool-mediated interaction between agents that represent different principals. PACT-PAIR measures the dyadic security--utility frontier. PACT-NET extends this to networked delegation, where third parties, cluster boundaries, and authority chains appear. The proposed solutions are relationship-specific policy, Mountable Context Cells, and pre-tool escalation. They are not final answers, but they identify promising control surfaces for safer coordination.


% ============================================================================
\section{What the Results Show}
\label{sec:conclusion_results}
% ============================================================================

The first result is that \textbf{privacy must be measured jointly with utility}. Refusal alone is not success. A delegate that refuses every request protects data but fails as a delegate. This is why the thesis uses the security--utility frontier rather than a single attack-success number. Across PACT-PAIR, the most useful policies are those that move both axes: they reduce disclosure while preserving legitimate task completion.

The second result is that \textbf{prompt specificity matters, but has a ceiling}. In PACT-PAIR, generic privacy instructions do not reliably move the frontier, while explicit category-level policies do. However, even the best prompt-level policies still leak under ambiguity, multi-turn interaction, and relationship variation. Multi-turn evaluation exposes channels that single-turn testing misses: facts can surface incidentally through related legitimate answers, refusal metadata, and accumulated context.

The third result is that \textbf{dyadic testing is not enough}. PACT-NET asks questions that PACT-PAIR cannot: whether information about a third party leaks through an intermediate agent, whether data crosses social or organisational clusters, whether one disclosure amplifies into multiple facts, and whether an agent accepts false delegation claims. These are not larger versions of the same task; they are network phenomena.

The fourth result is that \textbf{architectural control surfaces are promising but imperfect}. Relationship-specific policy gives the model a better boundary but still leaves protected data visible. MCC reduces leakage by removing unmounted context, but its effectiveness depends on scope granularity and separate read/write permissions. Escalation can gate risky tool calls before retrieval, but sparse precedent transfer remains representation-sensitive. The strongest supported conclusion is therefore not that any single mechanism solves privacy, but that useful governance requires a layered control plane.


% ============================================================================
\section{Practical Value}
\label{sec:practical_value}
% ============================================================================

The practical value of this work is a way to evaluate and build agent coordination systems without pretending that security and utility are separable. The target use cases are realistic: a founder's agent talking to an investor's agent, a recruiter coordinating with a candidate's agent, a team agent asking another team's agent for project context, or a personal agent interacting with an enterprise agent. In each case, the goal is not to prevent communication. The goal is to make useful communication possible without silently expanding access to everything the agent knows.

For builders, the thesis suggests a concrete design discipline. First, define the requester--owner relationship and the entitlement function before evaluating behaviour. Second, test both legitimate-task completion and protected-data disclosure. Third, separate read permissions from write permissions. Fourth, place enforcement before retrieval wherever possible: context scoping and tool gating are stronger than asking the model to ignore data it has already seen.

For research, SharedOS and PACT offer reusable infrastructure. New models, policies, social graphs, data surfaces, and attack strategies can be evaluated in the same setting. This matters because cross-boundary agent coordination will not be adopted if every new deployment relies on informal red-teaming or anecdotal trust. The field needs benchmarks that measure whether agents can help across boundaries while respecting those boundaries.


% ============================================================================
\section{Limitations}
\label{sec:limitations}
% ============================================================================

The results should be read with several limits in mind. The benchmark worlds are synthetic and centred on a technology-startup setting. Other domains, such as healthcare, legal work, or education, may produce different sensitivity boundaries and different failure distributions. The experiments primarily use same-ecosystem agents and do not exhaustively evaluate optimised adversarial attackers. The MCC implementation is strongest for note-folder scoping; todo isolation is currently coarser. The escalation experiment is an isolated pre-tool gate evaluation rather than a full end-to-end leakage run.

These limitations do not invalidate the core claim. They bound it. The thesis shows that cross-boundary delegation has measurable failure modes, that those modes change when interaction becomes multi-turn or networked, and that moving enforcement earlier in the stack can improve the frontier. It does not claim that the current mechanisms are complete.


% ============================================================================
\section{Conclusion}
\label{sec:closing}
% ============================================================================

Future agent systems will not consist of isolated assistants. They will consist of agents representing different people and organisations, negotiating access, sharing context, requesting actions, and escalating uncertainty. The hard problem is not only making each agent more capable. It is building the infrastructure that lets agents coordinate safely across trust boundaries.

This thesis argues that such infrastructure needs three properties. It must be \textbf{measurable}, because security without utility is not a useful target. It must be \textbf{relational}, because the same fact can be appropriate for one requester and private for another. It must be \textbf{structural}, because final-answer prompting is too late once protected data is already in context.

SharedOS, PACT-PAIR, PACT-NET, MCC, and escalation are steps toward that infrastructure. They frame cross-boundary agent privacy as a systems problem: define the boundary, measure the frontier, and move enforcement into the control plane. That framing is the main contribution. It makes safe multi-agent coordination a concrete engineering and research programme rather than a vague hope that agents will simply know what not to say.
